III) \emph{Fin de la demostración.} --- Estableceremos primero el lema
siguiente:

\begin{lemma}[6.14.1.2]
\label{IV.2.s6.14.1.2-fr}
Sean $A$ un anillo local noetheriano que es un dominio entero de
dimensión 1, $A'$ un anillo local noetheriano y $A\to A'$ un
homomorfismo local tal que el morfismo correspondiente
$\operatorname{Spec}(A')\to\operatorname{Spec}(A)$ sea normal. Sean $K$
el cuerpo de fracciones de $A$, $\mathfrak m$ el ideal maximal de $A$
y $k=A/\mathfrak m$ su cuerpo residual.
\oldpage[IV-2]{172}
\begin{enumerate}
\item[\textup{(i)}] Si $\mathfrak q'_j$ $(1\leqslant j\leqslant r)$ son los
ideales primos mínimos de $A'$, $K'=K\otimes_A A'$ es un producto directo
de dominios enteros integralmente cerrados $K'_j\supset A'/\mathfrak q'_j$ $(1\leqslant j\leqslant r)$, cada $K'_j$ con cuerpo de fracciones $L'_j$ igual al de
$A'/\mathfrak q'_j$ $(1\leqslant j\leqslant r)$, de modo que el anillo total de fracciones $L'$ de $A'$ se identifica con
el producto directo de los $L'_j$, y $A'$ con un subanillo de $K'$.
\item[\textup{(ii)}] El ideal $\mathfrak p'=\mathfrak mA'$ de $A'$ es primo; el
anillo $A'_{\mathfrak p'}$ tiene dimensión 1 y se identifica con un
subanillo del producto de los cuerpos $L'_j$ para los índices $j$
tales que $\mathfrak q'_j\subset\mathfrak p'$.
\item[\textup{(iii)}] Si $\overline{A}'_{\mathfrak p'}$ designa el subanillo de
$L'$ que es el producto de $A'_{\mathfrak p'}$ y de los $L'_j$ para
los que $\mathfrak q'_j\not\subset\mathfrak p'$, se tiene
\begin{equation}
\tag{6.14.1.3}
\label{IV.2.s6.14.1.3-fr}
A'=K'\cap\overline{A}'_{\mathfrak p'}.
\end{equation}
\end{enumerate}
\end{lemma}
